\documentclass[11pt]{article}

\usepackage[a4paper,margin=1in]{geometry}
\usepackage[T1]{fontenc}
\usepackage[utf8]{inputenc}
\usepackage{amsmath,amssymb,mathtools}
\usepackage{booktabs,longtable,array}
\usepackage{enumitem}
\usepackage{xcolor}
\usepackage{hyperref}
\usepackage{listings}

\hypersetup{
  colorlinks=true,
  linkcolor=blue!60!black,
  urlcolor=blue!60!black,
  citecolor=blue!60!black
}

\lstdefinestyle{pygloop}{
  basicstyle=\ttfamily\small,
  breaklines=true,
  columns=fullflexible,
  frame=single,
  framerule=0.2pt,
  backgroundcolor=\color{black!2}
}
\lstset{style=pygloop}

\newcommand{\dd}{\mathrm{d}}
\newcommand{\ii}{\mathrm{i}}
\newcommand{\GammaLoop}{\textsc{GammaLoop}}
\newcommand{\pygloop}{\texttt{pygloop}}
\newcommand{\qqbarNX}{\texttt{qqbar\_nX}}
\newcommand{\DOT}{\textsc{dot}}
\newcommand{\spat}[1]{\vec{#1}}
\newcommand{\Slash}[1]{\not\!#1}
\newcommand{\lzero}{\ell_0}
\newcommand{\lone}{\ell_1}
\newcommand{\D}{\mathcal{D}}
\newcommand{\N}{\mathcal{N}}
\newcommand{\I}{\mathcal{I}}
\newcommand{\MT}{m_t}
\newcommand{\rtab}[4]{#1 & #2 & \(#3\) & #4\\}

\title{\qqbarNX{} implementation notes: local two-loop initial-state collinear subtraction for \(d\bar d\to HHH\)}
\author{Generated implementation documentation}
\date{Updated 12 May 2026}

\begin{document}
\maketitle

\begin{abstract}
This document records the current implementation of the \qqbarNX{} process in
\pygloop{}, focusing on the phase-space symmetrized two-loop amplitude
\[
  d(p_0)\,\bar d(p_1)\to H(p_2)\,H(p_3)\,H(p_4).
\]
The selected diagrams contain a massive top-quark pentagon loop connected to
the incoming massless quark line by two gluons.  The local subtraction follows
the two initial-state collinear projectors described in
Ref.~\cite{KermanschahVicini2026}, but no UV subtraction is generated here:
the heavy subgraph is a pentagon and is UV finite for the present purpose.

The current \DOT{} encoding writes the original graphs and their local IR
counterterms as separate graphs sharing a \texttt{group\_id}.  Counterterm
numerators are stored globally as spenso tensor networks, while the CT
topology is used to expose the denominator structure to \GammaLoop{}'s CFF/LTD
machinery.  This note gives the exact four-dimensional representation of all
four selected symmetrized original graphs and all eight associated
counterterms, including the loop-momentum routing of every denominator.
\end{abstract}

\tableofcontents

\section{Current implementation state}

The relevant implementation files are:
\begin{description}[leftmargin=2.5cm,style=nextline]
\item[\texttt{config.toml}]
\path{src/processes/qqbar_nX/config.toml}.  User-facing process configuration.
\item[\texttt{qqbar\_nX.py}]
\path{src/processes/qqbar_nX/qqbar_nX.py}.  Process class, \GammaLoop{}
steering, run-card generation, and tests.
\item[\texttt{qqbar\_nX\_graphs.py}]
\path{src/processes/qqbar_nX/qqbar_nX_graphs.py}.  \DOT{} parsing, graph
selection, and canonical topology grouping.
\item[\texttt{qqbar\_nX\_counterterms.py}]
\path{src/processes/qqbar_nX/qqbar_nX_counterterms.py}.  Construction of ISR
counterterm \DOT{} graphs.
\item[\texttt{qqbar\_nX\_4d.py}]
\path{src/processes/qqbar_nX/qqbar_nX_4d.py}.  Direct four-dimensional
Symbolica/spenso evaluator.
\end{description}

The current default artifacts are:
\begin{description}[leftmargin=2.5cm,style=nextline]
\item[subtracted \DOT{}]
\path{outputs/dot_files/qqbar_nX/qqbar_nX_d_dbar_h_h_h_2L_top_pentagon_isr_subtracted.dot}.
\item[run card]
\path{outputs/dot_files/qqbar_nX/qqbar_nX_d_dbar_h_h_h_2L_top_pentagon_isr_subtracted.toml}.
\item[manifest]
\path{outputs/dot_files/qqbar_nX/qqbar_nX_d_dbar_h_h_h_2L_top_pentagon_isr_subtracted.manifest.json}.
\end{description}

The current default configuration is:
\begin{lstlisting}
[generation]
symmetrize_final_states = true

[counterterms]
build_isr = true
projector_mode = "leading"
denominator_strategy = "contract"
auxiliary_denominator_mode = "exact_xi_topology"
global_phase = "-i"
uv_inert_dod = -100
use_parametric_xi = false

[standalone_run_card]
enable_threshold_subtraction = true
\end{lstlisting}

For the symmetrized \(HHH\) final state the manifest records
\[
  N_{\rm original}=4,\qquad N_{\rm CT}=8,\qquad N_{\rm total}=12.
\]
The selected original graphs are \texttt{GL05}, \texttt{GL06},
\texttt{GL12}, and \texttt{GL13}.  Each original graph has two local
counterterms:
\[
  G_{\rm isr\_p1\_ct}\quad\hbox{for the \(p_0\)-collinear gluon bridge},
  \qquad
  G_{\rm isr\_p2\_ct}\quad\hbox{for the \(p_1\)-collinear gluon bridge}.
\]
The current groups are
\[
  \hbox{group 0}=\{\texttt{GL05},\texttt{GL06}\hbox{ and CTs}\},\qquad
  \hbox{group 1}=\{\texttt{GL12},\texttt{GL13}\hbox{ and CTs}\}.
\]
The original graph remains the group master; the CT graphs are marked
\texttt{is\_group\_master=false}.  All CT node and edge \texttt{dod}
attributes are forced to \(-100\), which prevents \GammaLoop{}'s own local UV
machinery from acting on these artificial local IR counterterms.

\section{Generation, selection, and grouping}

The raw generation command built by \texttt{qqbar\_nX.build\_feyngen\_command}
is conceptually
\begin{lstlisting}
generate amp d d~ > h h h | d d~ g h t t~ QED==3 [{2}]
  --only-diagrams
  --global-prefactor-projector '<external spin/color projector>'
  --numerator-grouping no_grouping
  --veto-vertex-interactions V_6 V_9 V_36 V_37
  --number-of-fermion-loops 1 1
  --graph-prefix GL
  --max-multiplicity-for-fast-cut-filter 99
  --symmetrize-final-states true
\end{lstlisting}
The selector in \texttt{qqbar\_nX\_graphs.py} retains a graph only if it has
exactly:
\begin{itemize}[leftmargin=2em]
  \item external particles \(d,\bar d,H,H,H\);
  \item two light-quark--gluon vertices \(\texttt{V\_74}\);
  \item two top--gluon vertices \(\texttt{V\_137}\);
  \item three top--Higgs vertices \(\texttt{V\_141}\);
  \item one internal massless light-quark edge joining the two light vertices;
  \item one gluon bridge from each light vertex to a top--gluon vertex;
  \item one connected top cycle with five top propagators.
\end{itemize}
Canonical topology grouping forgets raw node labels and uses label-preserving
permutations of the graph.  The resulting \texttt{group\_id} is the minimal
unit expected to be locally IR finite after adding the corresponding CTs.

\section{Momentum convention and loop-momentum basis}

The \DOT{} uses \GammaLoop{}'s all-incoming convention.  We write
\[
  p_0=P(0),\quad p_1=P(1),\quad p_2=P(2),\quad
  p_3=P(3),\quad p_4=P(4),
\]
and
\[
  \lzero=K(0),\qquad \lone=K(1).
\]
For every internal edge \(e\),
\[
  q_e^\mu
  =
  \rho_{e0}\ell_0^\mu+\rho_{e1}\ell_1^\mu
  +\sum_{r=0}^{4}\eta_{er}p_r^\mu ,
\]
as stored in the edge attribute \texttt{lmb\_rep}.  The four-dimensional
propagator denominator associated with edge \(e\) is
\[
  D_e=q_e^2-m_e^2+\ii0,\qquad
  m_e=
  \begin{cases}
    \MT, & e\hbox{ a top edge},\\
    0, & e\hbox{ a massless \(d\) or gluon edge}.
  \end{cases}
\]

The local subtraction depends on the routing of the light-line momentum.  The
implementation therefore pins the loop momentum basis explicitly:
\begin{center}
\begin{tabular}{ccll}
\toprule
group & originals & \texttt{lmb\_id=0} & \texttt{lmb\_id=1}\\
\midrule
0 & \texttt{GL05}, \texttt{GL06} & light bridge edge \(10\) & top edge \(5\)\\
1 & \texttt{GL12}, \texttt{GL13} & light bridge edge \(5\) & top edge \(8\)\\
\bottomrule
\end{tabular}
\end{center}
The same basis edges are preserved in every CT graph in the same group.  This
is important because \GammaLoop{} sets several group-level structures from the
group master.

\section{Exact four-dimensional integrand representation}

For each original or CT graph \(G\), the represented four-dimensional
integrand is
\begin{equation}
  \I_G(\ell_0,\ell_1)
  =
  \frac{\N_G(\ell_0,\ell_1)}
       {\prod_{e\in E_{\rm int}(G)}D_e(\ell_0,\ell_1)} ,
  \qquad
  A_G=\int\frac{\dd^4\ell_0}{(2\pi)^4}
          \frac{\dd^4\ell_1}{(2\pi)^4}\,
          \I_G(\ell_0,\ell_1).
  \label{eq:generalIntegrand}
\end{equation}
The denominator product is specified graph by graph in
Sec.~\ref{sec:denominatorTables}.

\subsection{Original graph numerators}

For original graphs,
\begin{equation}
  \N_G =
  {\cal P}_{\rm ext}
  \prod_{e\in E_{\rm ext}}E_{\rm ext}(e)
  \prod_{e\in E_{\rm int}}E_e(q_e)
  \prod_{v\in V_{\rm int}}V_v ,
  \label{eq:originalNumeratorProduct}
\end{equation}
with the graph-global external projector
\[
  {\cal P}_{\rm ext}
  =
  u(0)\,\bar v(1)\,\frac{1}{3}\,\delta_{\rm color}(\bar d,d).
\]
The local edge numerator rules exported by \GammaLoop{} are
\begin{align}
  E_d(e:a\to b)
  &=
  \ii\,q_{e,\mu}\,
  \gamma(s_b,s_a;\mu)\,
  \delta(c_a,c_b),
  \\
  E_t(e:a\to b)
  &=
  \ii\,[q_{e,\mu}\gamma(s_b,s_a;\mu)+\MT\,\delta(s_b,s_a)]
  \,\delta(c_a,c_b),
  \\
  E_g(e:a\to b)
  &=
  -\ii\,g(\mu_a,\mu_b)\,\delta(A_a,A_b),
  \\
  E_{\rm ext}(e)&=\ii
  \quad\hbox{for every external \(d,\bar d,H\) edge}.
\end{align}
The vertex numerator rules are
\begin{align}
  V_{74}(q q g)
  &=
  {\rm GC}_{11}\,\gamma(s_L,s_R;\mu_g)\,
  T(A_g,c_R,c_L),
  \\
  V_{137}(t t g)
  &=
  {\rm GC}_{11}\,\gamma(s_L,s_R;\mu_g)\,
  T(A_g,c_R,c_L),
  \\
  V_{141}(t t H)
  &=
  {\rm GC}_{94}\,[P_-(s_L,s_R)+P_+(s_L,s_R)]\,\delta(c_R,c_L).
\end{align}
In the independent four-dimensional evaluator \(P_-+P_+\) is simplified to
the spin identity.

\subsection{Counterterm graph numerators}

For CT graphs, the final \DOT{} has every local node and edge numerator set to
\(\texttt{num}=1\).  The complete numerator is collected first and then stored
in the graph-global \texttt{num} attribute as a spenso tensor network.  In
factorised notation the current default CT numerator is
\begin{equation}
  \N_{{\rm CT}}(G,b)
  =
  \Phi_{\rm CFF}\,\eta_b\,R_b(G)\,A_b\,B_b(G)^{-1}\,
  {\cal C}_{G,b}[\N_G],
  \qquad b=1,2.
  \label{eq:ctNumerator}
\end{equation}
The current default is
\[
  \Phi_{\rm CFF}=-\ii,\qquad
  \eta_1=-\ii,\qquad
  \eta_2=+\ii.
\]
The additional phase \(\Phi_{\rm CFF}\) is a current \GammaLoop{} CFF/event
convention phase.  It is required for the 3D CFF inspect cancellation.  For a
pure paper-level four-dimensional Feynman-integrand identity, this phase is
set to \(1\) in a temporary config.

The fixed auxiliary vector is
\[
  \zeta=p_0+p_1,
\]
and the default implementation uses \texttt{use\_parametric\_xi=false}.  The
Lorentz prefactors are
\[
  A_1=\frac{p_0\cdot\zeta}{p_0\cdot p_1},
  \qquad
  A_2=\frac{p_1\cdot\zeta}{p_0\cdot p_1}.
\]
For massless beams and \(\zeta=p_0+p_1\), both evaluate to one, but the DOT
keeps the tensor expression.

Define the Euclidean spatial projection
\[
  \langle a,b\rangle_3=a_xb_x+a_yb_y+a_zb_z,\qquad
  |a|_3^2=\langle a,a\rangle_3.
\]
Let \(L\) be the pinned light bridge, \(g_1\) the gluon that is collinear to
\(p_0\), and \(g_2\) the gluon that is collinear to \(p_1\).  The finite
routing fractions are
\begin{equation}
  R_1(G)=\frac{\langle q_L,p_1\rangle_3}{\langle q_{g_1},p_1\rangle_3},
  \qquad
  R_2(G)=\frac{\langle q_L,p_0\rangle_3}{\langle q_{g_2},p_0\rangle_3}.
  \label{eq:routingFractions}
\end{equation}
They deliberately use only spatial components and therefore introduce no
loop-energy poles in the graph-global numerator.  This was a diagnostic
implementation used before the exact auxiliary topology was available; it is
not the default subtraction encoded in the current \DOT{}.

The current CFF-compatible auxiliary denominator mode is
\texttt{exact\_xi\_topology}.  Following the author-confirmed correction of
Eqs.~(12)--(13) in Ref.~\cite{KermanschahVicini2026}, the auxiliary
propagator is built from the adjacent gluon momentum \(k_{g_i}\), not from the
light-bridge momentum \(k_1\).  Define
\begin{equation}
  A_\xi(q)\equiv (q-\xi)^2-\xi^2=q^2-2q\cdot\xi .
  \label{eq:correctedAuxiliary}
\end{equation}
The light-line denominator structures are therefore
\begin{align}
  D_{\rm orig} &= k_1^2\,k_{g1}^2\,k_{g2}^2,\\
  D_{\Delta_1} &= k_1^2\,k_{g1}^2\,A_\xi(k_{g1}),\\
  D_{\Delta_2} &= k_1^2\,A_\xi(k_{g2})\,k_{g2}^2 .
  \label{eq:correctedDenominators}
\end{align}
In \DOT{} form, this is realized by splitting the relevant collinear gluon
edge.  The original gluon edge id is kept on the unshifted segment so
\(k_{g_i}^2\) remains an ordinary topology denominator.  A new massive segment
is inserted on the same path and the active fake in/out helper pair is attached
around it; when the helpers are sampled as \(Q(2)=Q(4)=\xi\) for \(\Delta_1\)
or \(Q(3)=Q(5)=\xi\) for \(\Delta_2\), GammaLoop routes this segment as
\(k_{g_i}-\xi\).  The segment mass is the corresponding helper invariant
\(\xi^2\), so the topology denominator is exactly \(A_\xi(k_{g_i})\).  The
opposite gluon denominator is removed by the configured contraction strategy.
No ratio or inverse auxiliary denominator is placed in the graph-global
numerator in this mode.

The legacy \texttt{spatial\_proxy} mode removes the opposite gluon propagator
from the CT topology and keeps the leading fixed-\(\zeta\) auxiliary
denominator as a loop-energy-free graph numerator factor:
\begin{align}
  B_1(G)
  &=
  -2(p_0\cdot p_1)
  \frac{\langle q_L,p_0\rangle_3}{|p_0|_3^2},
  \\
  B_2(G)
  &=
  +2(p_0\cdot p_1)
  \frac{\langle q_L,p_0\rangle_3}{|p_0|_3^2}.
  \label{eq:spatialProxy}
\end{align}
In the current exact topology this expression is no longer part of the default
graph numerator; it is retained only as a diagnostic of the old implementation.
Denominators with actual loop-energy dependence are part of the graph topology.

The transformation \({\cal C}_{G,b}\) acts on the original numerator as:
\begin{itemize}[leftmargin=2em]
  \item replace the internal light-bridge numerator by
        \(\delta_{\rm spin}\delta_{\rm color}\);
  \item for \(b=1\), replace the \(p_0\)-side light vertex by
        \[
          V_{{\rm CT},1}
          =
          \omega_1\,{\rm GC}_{11}\,
          T(A_{g_1},c_{\rm ext},c_{\rm int})\,
          [2q_{g_1}^{\mu_{g_1}}\delta_{\rm spin}];
        \]
  \item for \(b=2\), replace the \(p_1\)-side light vertex by
        \[
          V_{{\rm CT},2}
          =
          \omega_2\,{\rm GC}_{11}\,
          T(A_{g_2},c_{\rm int},c_{\rm ext})\,
          [-2q_{g_2}^{\mu_{g_2}}\delta_{\rm spin}];
        \]
  \item split the same-side collinear gluon edge to introduce the topology
        denominator \(A_\xi(k_{g_i})\);
  \item contract the opposite gluon edge into the opposite light endpoint,
        multiplying the opposite light vertex, the opposite gluon numerator,
        and the corresponding top--gluon vertex into the graph-global
        numerator.
\end{itemize}
The orientation signs currently selected from the \DOT{} are:
\begin{center}
\begin{tabular}{lcc}
\toprule
original graph & \(\omega_1\) & \(\omega_2\)\\
\midrule
\texttt{GL05}, \texttt{GL06} & \(+1\) & \(-1\)\\
\texttt{GL12}, \texttt{GL13} & \(-1\) & \(+1\)\\
\bottomrule
\end{tabular}
\end{center}

\subsection[Relation to reference subtraction formula]{Relation to Ref.~\cite{KermanschahVicini2026}}

The paper writes the light-line factor schematically as
\[
  D_{\alpha\beta}
  =
  \ii\,
  \frac{
    \bar v(p_1)\gamma_\beta(-\Slash{k}_1)\gamma_\alpha u(p_0)
  }{
    k_1^2\,k_{g1}^2\,k_{g2}^2
  },
\]
and constructs two local operators \(\Delta_1,\Delta_2\) that replace the
light-line numerator by the leading collinear numerator and introduce an
auxiliary denominator built from a fixed vector.  In the published text that
auxiliary denominator is written with \(k_1\); the author has confirmed this is
a typo and that the intended momenta are \(k_{g1}\) and \(k_{g2}\), as encoded
in Eq.~\eqref{eq:correctedDenominators}.  The \qqbarNX{} code uses the same
local leading-collinear replacement logic for the two ISR regions, but differs
in two implementation-level ways:
\begin{enumerate}[leftmargin=2em]
  \item the heavy subgraph is a top pentagon with three Higgs insertions;
  \item no UV subtraction operator is generated.
\end{enumerate}

\section{Original numerator topologies}
\label{sec:numeratorTopologies}

The original numerator of each selected graph is exactly obtained by applying
Eq.~\eqref{eq:originalNumeratorProduct} and the local rules above to the
following directed \DOT{} topologies.  An entry \(u:h\to v:k\) denotes source
node \(u\) at hedge \(h\) and destination node \(v\) at hedge \(k\).

\begin{longtable}{>{\ttfamily}p{0.14\linewidth}p{0.25\linewidth}p{0.50\linewidth}}
\toprule
graph & internal nodes & internal edges \\
\midrule
\endhead
GL05 &
0 V\_137, 1 V\_137, 2 V\_74, 3 V\_74, 4 V\_141, 5 V\_141, 6 V\_141 &
e5 t \(0{:}5\to1{:}6\);
e6 g \(0{:}7\to3{:}8\);
e7 t \(6{:}9\to0{:}10\);
e8 g \(1{:}11\to2{:}12\);
e9 t \(1{:}13\to5{:}14\);
e10 d \(3{:}15\to2{:}16\);
e11 t \(5{:}17\to4{:}18\);
e12 t \(4{:}19\to6{:}20\).\\
\addlinespace
GL06 &
0 V\_137, 1 V\_137, 2 V\_74, 3 V\_74, 4 V\_141, 5 V\_141, 6 V\_141 &
e5 t \(1{:}5\to0{:}6\);
e6 g \(0{:}7\to3{:}8\);
e7 t \(0{:}9\to6{:}10\);
e8 g \(1{:}11\to2{:}12\);
e9 t \(5{:}13\to1{:}14\);
e10 d \(3{:}15\to2{:}16\);
e11 t \(4{:}17\to5{:}18\);
e12 t \(6{:}19\to4{:}20\).\\
\addlinespace
GL12 &
0 V\_74, 1 V\_74, 2 V\_137, 3 V\_137, 4 V\_141, 5 V\_141, 6 V\_141 &
e5 d \(1{:}5\to0{:}6\);
e6 g \(0{:}7\to3{:}8\);
e7 g \(1{:}9\to2{:}10\);
e8 t \(2{:}11\to5{:}12\);
e9 t \(6{:}13\to2{:}14\);
e10 t \(4{:}15\to3{:}16\);
e11 t \(3{:}17\to6{:}18\);
e12 t \(5{:}19\to4{:}20\).\\
\addlinespace
GL13 &
0 V\_74, 1 V\_74, 2 V\_137, 3 V\_137, 4 V\_141, 5 V\_141, 6 V\_141 &
e5 d \(1{:}5\to0{:}6\);
e6 g \(0{:}7\to3{:}8\);
e7 g \(1{:}9\to2{:}10\);
e8 t \(5{:}11\to2{:}12\);
e9 t \(2{:}13\to6{:}14\);
e10 t \(3{:}15\to4{:}16\);
e11 t \(6{:}17\to3{:}18\);
e12 t \(4{:}19\to5{:}20\).\\
\bottomrule
\end{longtable}

\section{Exact denominator routings}
\label{sec:denominatorTables}

This section gives the exact four-dimensional denominator routing for every
graph in the current symmetrized subtracted \DOT{}.  For each listed graph,
\[
  \I_G=\frac{\N_G}{\prod_{e\in S_G}D_e},
\]
where \(S_G\) is the edge set shown in the corresponding product.  For CT
graphs, \(\N_G\) means the global numerator of Eq.~\eqref{eq:ctNumerator} with
the original graph and beam indicated by the graph name.

\subsection{\texttt{GL05} and counterterms}

\paragraph{\texttt{GL05}.}
\[
  \I_{\texttt{GL05}}=
  \frac{\N_{\texttt{GL05}}}{D_5D_6D_7D_8D_9D_{10}D_{11}D_{12}}.
\]
\begin{center}
\begin{tabular}{ccll}
\toprule
edge & part. & \(q_e\) & basis\\
\midrule
\rtab{5}{t}{\lzero}{\texttt{lmb\_id=1}}
\rtab{6}{g}{-\lzero-p_0-p_1+\lone+p_3}{}
\rtab{7}{t}{-p_0-p_1+\lone+p_3}{}
\rtab{8}{g}{\lzero-\lone-p_3}{}
\rtab{9}{t}{\lone+p_3}{}
\rtab{10}{d}{-\lzero-p_1+\lone+p_3}{\texttt{lmb\_id=0}}
\rtab{11}{t}{\lone}{}
\rtab{12}{t}{-p_0-p_1+\lone+p_2+p_3}{}
\bottomrule
\end{tabular}
\end{center}

\paragraph{\texttt{GL05\_isr\_p1\_ct}.}
This is \(N_{\rm CT}(\texttt{GL05},1)\).  It keeps \(g_1=e6\) and removes
the opposite gluon \(e8\) from the active denominator topology:
\[
  \I_{\texttt{GL05\_isr\_p1\_ct}}=
  \frac{\N_{\rm CT}(\texttt{GL05},1)}
       {D_5D_6D_7D_8D_9D_{10}D_{11}}.
\]
\begin{center}
\begin{tabular}{ccll}
\toprule
edge & part. & \(q_e\) & basis\\
\midrule
\rtab{5}{t}{\lzero}{\texttt{lmb\_id=1}}
\rtab{6}{g}{-\lzero-p_0-p_1+\lone+p_3}{}
\rtab{7}{t}{-p_0-p_1+\lone+p_3}{}
\rtab{8}{t}{-p_0-p_1+\lone+p_2+p_3}{}
\rtab{9}{t}{\lone+p_3}{}
\rtab{10}{d}{-\lzero-p_1+\lone+p_3}{\texttt{lmb\_id=0}}
\rtab{11}{t}{\lone}{}
\bottomrule
\end{tabular}
\end{center}

\paragraph{\texttt{GL05\_isr\_p2\_ct}.}
This is \(N_{\rm CT}(\texttt{GL05},2)\).  It keeps \(g_2=e8\) and removes
the opposite gluon \(e6\):
\[
  \I_{\texttt{GL05\_isr\_p2\_ct}}=
  \frac{\N_{\rm CT}(\texttt{GL05},2)}
       {D_5D_6D_7D_8D_9D_{10}D_{11}}.
\]
\begin{center}
\begin{tabular}{ccll}
\toprule
edge & part. & \(q_e\) & basis\\
\midrule
\rtab{5}{t}{\lzero}{\texttt{lmb\_id=1}}
\rtab{6}{t}{-p_0-p_1+\lone+p_2+p_3}{}
\rtab{7}{t}{-p_0-p_1+\lone+p_3}{}
\rtab{8}{g}{\lzero-\lone-p_3}{}
\rtab{9}{t}{\lone+p_3}{}
\rtab{10}{d}{-\lzero-p_1+\lone+p_3}{\texttt{lmb\_id=0}}
\rtab{11}{t}{\lone}{}
\bottomrule
\end{tabular}
\end{center}

\subsection{\texttt{GL06} and counterterms}

\paragraph{\texttt{GL06}.}
\[
  \I_{\texttt{GL06}}=
  \frac{\N_{\texttt{GL06}}}{D_5D_6D_7D_8D_9D_{10}D_{11}D_{12}}.
\]
\begin{center}
\begin{tabular}{ccll}
\toprule
edge & part. & \(q_e\) & basis\\
\midrule
\rtab{5}{t}{\lzero}{\texttt{lmb\_id=1}}
\rtab{6}{g}{\lzero-\lone-p_0-p_1+p_3}{}
\rtab{7}{t}{\lone+p_0+p_1-p_3}{}
\rtab{8}{g}{-\lzero+\lone-p_3}{}
\rtab{9}{t}{\lone-p_3}{}
\rtab{10}{d}{\lzero-\lone-p_1+p_3}{\texttt{lmb\_id=0}}
\rtab{11}{t}{\lone}{}
\rtab{12}{t}{\lone+p_0+p_1-p_2-p_3}{}
\bottomrule
\end{tabular}
\end{center}

\paragraph{\texttt{GL06\_isr\_p1\_ct}.}
\[
  \I_{\texttt{GL06\_isr\_p1\_ct}}=
  \frac{\N_{\rm CT}(\texttt{GL06},1)}
       {D_5D_6D_7D_8D_9D_{10}D_{11}}.
\]
\begin{center}
\begin{tabular}{ccll}
\toprule
edge & part. & \(q_e\) & basis\\
\midrule
\rtab{5}{t}{\lzero}{\texttt{lmb\_id=1}}
\rtab{6}{g}{\lzero-\lone-p_0-p_1+p_3}{}
\rtab{7}{t}{\lone+p_0+p_1-p_3}{}
\rtab{8}{t}{\lone+p_0+p_1-p_2-p_3}{}
\rtab{9}{t}{\lone-p_3}{}
\rtab{10}{d}{\lzero-\lone-p_1+p_3}{\texttt{lmb\_id=0}}
\rtab{11}{t}{\lone}{}
\bottomrule
\end{tabular}
\end{center}

\paragraph{\texttt{GL06\_isr\_p2\_ct}.}
\[
  \I_{\texttt{GL06\_isr\_p2\_ct}}=
  \frac{\N_{\rm CT}(\texttt{GL06},2)}
       {D_5D_6D_7D_8D_9D_{10}D_{11}}.
\]
\begin{center}
\begin{tabular}{ccll}
\toprule
edge & part. & \(q_e\) & basis\\
\midrule
\rtab{5}{t}{\lzero}{\texttt{lmb\_id=1}}
\rtab{6}{t}{\lone+p_0+p_1-p_2-p_3}{}
\rtab{7}{t}{\lone+p_0+p_1-p_3}{}
\rtab{8}{g}{-\lzero+\lone-p_3}{}
\rtab{9}{t}{\lone-p_3}{}
\rtab{10}{d}{\lzero-\lone-p_1+p_3}{\texttt{lmb\_id=0}}
\rtab{11}{t}{\lone}{}
\bottomrule
\end{tabular}
\end{center}

\subsection{\texttt{GL12} and counterterms}

\paragraph{\texttt{GL12}.}
\[
  \I_{\texttt{GL12}}=
  \frac{\N_{\texttt{GL12}}}{D_5D_6D_7D_8D_9D_{10}D_{11}D_{12}}.
\]
\begin{center}
\begin{tabular}{ccll}
\toprule
edge & part. & \(q_e\) & basis\\
\midrule
\rtab{5}{d}{-\lzero+p_0}{\texttt{lmb\_id=0}}
\rtab{6}{g}{-\lzero+p_0+p_1}{}
\rtab{7}{g}{\lzero}{}
\rtab{8}{t}{\lone+p_0+p_1-p_2}{\texttt{lmb\_id=1}}
\rtab{9}{t}{-\lzero+\lone+p_0+p_1-p_2}{}
\rtab{10}{t}{\lone}{}
\rtab{11}{t}{-\lzero+\lone+p_0+p_1}{}
\rtab{12}{t}{\lone+p_0+p_1-p_2-p_3}{}
\bottomrule
\end{tabular}
\end{center}

\paragraph{\texttt{GL12\_isr\_p1\_ct}.}
This keeps \(g_1=e7\) and removes \(e6\):
\[
  \I_{\texttt{GL12\_isr\_p1\_ct}}=
  \frac{\N_{\rm CT}(\texttt{GL12},1)}
       {D_5D_6D_7D_8D_9D_{10}D_{11}}.
\]
\begin{center}
\begin{tabular}{ccll}
\toprule
edge & part. & \(q_e\) & basis\\
\midrule
\rtab{5}{d}{-\lzero+p_0}{\texttt{lmb\_id=0}}
\rtab{6}{t}{\lone+p_0+p_1-p_2-p_3}{}
\rtab{7}{g}{\lzero}{}
\rtab{8}{t}{\lone+p_0+p_1-p_2}{\texttt{lmb\_id=1}}
\rtab{9}{t}{-\lzero+\lone+p_0+p_1-p_2}{}
\rtab{10}{t}{\lone}{}
\rtab{11}{t}{-\lzero+\lone+p_0+p_1}{}
\bottomrule
\end{tabular}
\end{center}

\paragraph{\texttt{GL12\_isr\_p2\_ct}.}
This keeps \(g_2=e6\) and removes \(e7\):
\[
  \I_{\texttt{GL12\_isr\_p2\_ct}}=
  \frac{\N_{\rm CT}(\texttt{GL12},2)}
       {D_5D_6D_7D_8D_9D_{10}D_{11}}.
\]
\begin{center}
\begin{tabular}{ccll}
\toprule
edge & part. & \(q_e\) & basis\\
\midrule
\rtab{5}{d}{-\lzero+p_0}{\texttt{lmb\_id=0}}
\rtab{6}{g}{-\lzero+p_0+p_1}{}
\rtab{7}{t}{\lone+p_0+p_1-p_2-p_3}{}
\rtab{8}{t}{\lone+p_0+p_1-p_2}{\texttt{lmb\_id=1}}
\rtab{9}{t}{-\lzero+\lone+p_0+p_1-p_2}{}
\rtab{10}{t}{\lone}{}
\rtab{11}{t}{-\lzero+\lone+p_0+p_1}{}
\bottomrule
\end{tabular}
\end{center}

\subsection{\texttt{GL13} and counterterms}

\paragraph{\texttt{GL13}.}
\[
  \I_{\texttt{GL13}}=
  \frac{\N_{\texttt{GL13}}}{D_5D_6D_7D_8D_9D_{10}D_{11}D_{12}}.
\]
\begin{center}
\begin{tabular}{ccll}
\toprule
edge & part. & \(q_e\) & basis\\
\midrule
\rtab{5}{d}{-\lzero+p_0}{\texttt{lmb\_id=0}}
\rtab{6}{g}{-\lzero+p_0+p_1}{}
\rtab{7}{g}{\lzero}{}
\rtab{8}{t}{\lone-p_0-p_1+p_2}{\texttt{lmb\_id=1}}
\rtab{9}{t}{\lzero+\lone-p_0-p_1+p_2}{}
\rtab{10}{t}{\lone}{}
\rtab{11}{t}{\lzero+\lone-p_0-p_1}{}
\rtab{12}{t}{\lone-p_0-p_1+p_2+p_3}{}
\bottomrule
\end{tabular}
\end{center}

\paragraph{\texttt{GL13\_isr\_p1\_ct}.}
This keeps \(g_1=e7\) and removes \(e6\):
\[
  \I_{\texttt{GL13\_isr\_p1\_ct}}=
  \frac{\N_{\rm CT}(\texttt{GL13},1)}
       {D_5D_6D_7D_8D_9D_{10}D_{11}}.
\]
\begin{center}
\begin{tabular}{ccll}
\toprule
edge & part. & \(q_e\) & basis\\
\midrule
\rtab{5}{d}{-\lzero+p_0}{\texttt{lmb\_id=0}}
\rtab{6}{t}{\lone-p_0-p_1+p_2+p_3}{}
\rtab{7}{g}{\lzero}{}
\rtab{8}{t}{\lone-p_0-p_1+p_2}{\texttt{lmb\_id=1}}
\rtab{9}{t}{\lzero+\lone-p_0-p_1+p_2}{}
\rtab{10}{t}{\lone}{}
\rtab{11}{t}{\lzero+\lone-p_0-p_1}{}
\bottomrule
\end{tabular}
\end{center}

\paragraph{\texttt{GL13\_isr\_p2\_ct}.}
This keeps \(g_2=e6\) and removes \(e7\):
\[
  \I_{\texttt{GL13\_isr\_p2\_ct}}=
  \frac{\N_{\rm CT}(\texttt{GL13},2)}
       {D_5D_6D_7D_8D_9D_{10}D_{11}}.
\]
\begin{center}
\begin{tabular}{ccll}
\toprule
edge & part. & \(q_e\) & basis\\
\midrule
\rtab{5}{d}{-\lzero+p_0}{\texttt{lmb\_id=0}}
\rtab{6}{g}{-\lzero+p_0+p_1}{}
\rtab{7}{t}{\lone-p_0-p_1+p_2+p_3}{}
\rtab{8}{t}{\lone-p_0-p_1+p_2}{\texttt{lmb\_id=1}}
\rtab{9}{t}{\lzero+\lone-p_0-p_1+p_2}{}
\rtab{10}{t}{\lone}{}
\rtab{11}{t}{\lzero+\lone-p_0-p_1}{}
\bottomrule
\end{tabular}
\end{center}

\section{Four-dimensional evaluator}

The independent 4D test path in \texttt{qqbar\_nX\_4d.py} constructs
\[
  \I_{\mathcal G}^{4D}
  =
  \frac{
    {\rm Contr}\!\left[
      {\cal P}_{\rm ext}\,
      {\cal N}_{\mathcal G}^{\rm graph}
      \prod_v{\cal N}_v
      \prod_e{\cal N}_e
    \right]
  }{
    \prod_{e\in E_{\rm int}\setminus E_{\rm dummy}}
    \left(Q_e(0)^2-Q_e(1)^2-Q_e(2)^2-Q_e(3)^2-m_e^2\right)
  }.
  \label{eq:fourDEvaluator}
\]
It replaces \(\texttt{spenso::projm}+\texttt{spenso::projp}\) by the Dirac
metric, simplifies color with Symbolica community, and contracts the tensor
network with spenso before numerical evaluation.  If a graph-global numerator
contains \({\rm OSE}(e)\), it is replaced in the 4D diagnostic by
\[
  {\rm OSE}(e)\to
  \sqrt{Q_e(1)^2+Q_e(2)^2+Q_e(3)^2+m_e^2}.
\]

Implementation caveat: the current DOTs do not write an explicit top mass
attribute on top edges.  The formal \GammaLoop{} interpretation uses
\(m_t={\rm MT}\) for top denominators; the current independent 4D evaluator
only subtracts masses present as DOT edge attributes.  The tested collinear
identity is nevertheless meaningful because the heavy-loop denominators are
finite in the collinear approach and the same DOT object is evaluated for the
originals and their CTs.  A later cleanup should infer model masses by
particle name.

\section{Collinear tests and current results}

The process-specific tests are exposed through:
\begin{lstlisting}
PYTHONPATH=src .venv/bin/python src/pygloop.py --process qqbar_nX test --mode 4D
PYTHONPATH=src .venv/bin/python src/pygloop.py --process qqbar_nX test --mode gammaloop
\end{lstlisting}
The current test configuration uses \(x=0.3\), precision
\texttt{ArbPrec}, and
\[
  |k_\perp|=\lambda\sqrt{s},
  \qquad
  \lambda\in\{1,10^{-1},10^{-2},10^{-3},10^{-4},10^{-5},10^{-6},10^{-7}\}.
\]
For every \(\lambda\) and group, the diagnostic ratio is
\[
  \rho_{\rm group}
  =
  \frac{\left|\sum_{i\in{\rm group}}\I_i\right|}
       {\sum_{i\in{\rm group}}|\I_i|}.
\]
A local cancellation of the leading collinear singularity appears as a ratio
that decreases linearly with \(\lambda\).

\subsection[Pure 4D paper-level diagnostic, Phi CFF = 1]{Pure 4D paper-level diagnostic, \(\Phi_{\rm CFF}=1\)}

The independent 4D algebra check should be run with the CT global phase reset
to \(1\), because it tests the paper-level Feynman-integrand identity, not the
\GammaLoop{} 3D CFF event convention:
\begin{lstlisting}
tmp=/tmp/qqbar_nX_4d_phase1.toml
cp src/processes/qqbar_nX/config.toml "$tmp"
perl -0pi -e 's/global_phase = "-i"/global_phase = "1"/' "$tmp"
PYTHONPATH=src .venv/bin/python src/pygloop.py --process qqbar_nX \
  --qqbar-nx-config "$tmp" \
  --overwrite-process-basename qqbar_nX_4dphase1 \
  --m_top 1500 test --mode 4D
\end{lstlisting}
The current results are:
\begin{center}
\begin{tabular}{ccccc}
\toprule
beam & group & \(\lambda=10^{-3}\) & \(\lambda=10^{-5}\) & \(\lambda=10^{-7}\)\\
\midrule
\(p_0\) & 0 & \(1.207432422884\times10^{-2}\) & \(1.219429826526\times10^{-4}\) & \(1.219547309408\times10^{-6}\)\\
\(p_0\) & 1 & \(5.953650138218\times10^{-3}\) & \(5.918551126568\times10^{-5}\) & \(5.918194450500\times10^{-7}\)\\
\(p_1\) & 0 & \(5.110458964993\times10^{-3}\) & \(5.133079576218\times10^{-5}\) & \(5.133298939401\times10^{-7}\)\\
\(p_1\) & 1 & \(5.329260932463\times10^{-3}\) & \(5.304017780889\times10^{-5}\) & \(5.303763375967\times10^{-7}\)\\
\bottomrule
\end{tabular}
\end{center}
This confirms that the 4D numerator/projector construction cancels locally.

\subsection[Default 4D diagnostic with Phi CFF = -i]{Default 4D diagnostic with \(\Phi_{\rm CFF}=-\ii\)}

With the default CFF phase, the same 4D Feynman-integrand diagnostic does not
cancel, as expected:
\begin{center}
\begin{tabular}{ccc}
\toprule
beam & group & \(\rho_{\rm group}(\lambda=10^{-7})\)\\
\midrule
\(p_0\) & 0 & \(6.932501873977\times10^{-1}\)\\
\(p_0\) & 1 & \(7.070605995487\times10^{-1}\)\\
\(p_1\) & 0 & \(5.417365569681\times10^{-1}\)\\
\(p_1\) & 1 & \(7.068718138865\times10^{-1}\)\\
\bottomrule
\end{tabular}
\end{center}
This is a convention check: the \(-\ii\) phase is inserted for the
\GammaLoop{} CFF representation, not for the pure 4D algebra identity.

\subsection[GammaLoop 3D CFF inspect, mt = 1500]{\GammaLoop{} 3D CFF inspect, \(m_t=1500\)}

The non-threshold-clean heavy-mass test was run with:
\begin{lstlisting}
PYTHONPATH=src .venv/bin/python src/pygloop.py --process qqbar_nX \
  --m_top 1500 --clean test --mode gammaloop
\end{lstlisting}
The current results are:
\begin{center}
\begin{tabular}{ccccc}
\toprule
beam & group & \(\lambda=10^{-3}\) & \(\lambda=10^{-5}\) & \(\lambda=10^{-7}\)\\
\midrule
\(p_0\) & 0 & \(2.156437298117\times10^{-3}\) & \(2.162682427278\times10^{-5}\) & \(2.162745021408\times10^{-7}\)\\
\(p_0\) & 1 & \(1.019780744626\times10^{-2}\) & \(1.015716257667\times10^{-4}\) & \(1.015675267491\times10^{-6}\)\\
\(p_1\) & 0 & \(2.156437298117\times10^{-3}\) & \(2.162682427278\times10^{-5}\) & \(2.162745021408\times10^{-7}\)\\
\(p_1\) & 1 & \(1.019780744626\times10^{-2}\) & \(1.015716257667\times10^{-4}\) & \(1.015675267491\times10^{-6}\)\\
\bottomrule
\end{tabular}
\end{center}
The \(p_0\) and \(p_1\) summaries are identical in the current fractional
test because the GammaLoop inspect path solves the same pinned light-edge ray.
The important result is the linear decrease with \(\lambda\) in each group.

\subsection[GammaLoop 3D CFF inspect, physical mt = 173]{\GammaLoop{} 3D CFF inspect, physical \(m_t=173\)}

The physical-mass test was run at \(\sqrt{s}=1000\) GeV with threshold
subtraction enabled:
\begin{lstlisting}
PYTHONPATH=src .venv/bin/python src/pygloop.py --process qqbar_nX \
  --overwrite-process-basename qqbar_nX_mt173 \
  --m_top 173 --clean test --mode gammaloop
\end{lstlisting}
The generated threshold overlap structures were active:
\[
  \hbox{group 0 sizes }[13,12,12],
  \qquad
  \hbox{group 1 sizes }[9,9,8,8].
\]
The current results are:
\begin{center}
\begin{tabular}{ccccc}
\toprule
beam & group & \(\lambda=10^{-3}\) & \(\lambda=10^{-5}\) & \(\lambda=10^{-7}\)\\
\midrule
\(p_0\) & 0 & \(1.388028023261\times10^{-3}\) & \(1.409386585386\times10^{-5}\) & \(1.409618038701\times10^{-7}\)\\
\(p_0\) & 1 & \(1.738779562355\times10^{-3}\) & \(1.729480510655\times10^{-5}\) & \(1.729359476107\times10^{-7}\)\\
\(p_1\) & 0 & \(1.388028023261\times10^{-3}\) & \(1.409386585386\times10^{-5}\) & \(1.409618038701\times10^{-7}\)\\
\(p_1\) & 1 & \(1.738779562355\times10^{-3}\) & \(1.729480510655\times10^{-5}\) & \(1.729359476107\times10^{-7}\)\\
\bottomrule
\end{tabular}
\end{center}
This demonstrates the same local cancellation pattern in the physical-mass
setup with threshold sectors present.

\section{Standalone \GammaLoop{} run card}

The generated run card has named command blocks and no non-empty top-level
\texttt{commands} block, so loading the TOML does not automatically run
generation or inspection.  The intended standalone sequence is:
\begin{lstlisting}
GL=/Users/vjhirsch/Documents/Work/gammaloop_dev/target/dev-optim/gammaloop
CARD=outputs/dot_files/qqbar_nX/qqbar_nX_d_dbar_h_h_h_2L_top_pentagon_isr_subtracted.toml
STATE=outputs/gammaloop_states/qqbar_nX_standalone_qqbar_nX_d_dbar_h_h_h_2L_top_pentagon_isr_subtracted

$GL --clean-state "$CARD"
$GL -o -s "$STATE" run generate_subtracted_integrand
$GL -o -s "$STATE" run inspect_collinear_p1
$GL -o -s "$STATE" run inspect_collinear_p1_arb
$GL -o -s "$STATE" run inspect_collinear_p2
$GL -o -s "$STATE" run inspect_collinear_p2_arb
$GL -o -s "$STATE" run low_stat_integrate_pm
$GL -o -s "$STATE" run low_stat_integrate_pp
\end{lstlisting}
The helper shell script supports sectioned execution, including an \texttt{open}
mode that drops into the live \GammaLoop{} CLI.

\section{Current caveats and next cleanup items}

\begin{enumerate}[leftmargin=2em]
  \item The independent 4D evaluator should infer masses from the model or
        particle name, not only from explicit DOT edge \texttt{mass}
        attributes.
  \item The parametric-\(\xi\) option is implemented as configuration and
        additional-parameter plumbing, but the current validated results use
        the fixed \(\zeta=p_0+p_1\) setup.
  \item The current GammaLoop inspect \(p_0\) and \(p_1\) summaries use the
        same pinned light-edge fractional ray.  This is a useful grouped
        regression test, but a later diagnostic should also print independent
        solved rays for the two beam limits.
\end{enumerate}

\begin{thebibliography}{9}
\bibitem{KermanschahVicini2026}
D.~Kermanschah and M.~Vicini,
\emph{Heavy-quark box-loop corrections to \(q\bar q \to Z\gamma\) at two loops in QCD},
arXiv:2603.03169 [hep-ph],
\url{https://arxiv.org/abs/2603.03169}.
\end{thebibliography}

\end{document}
